\begin{ThreePartTable}
\begin{TableNotes}
    \item[1] The Xtensa processor has 64 internal AR registers. It is designed with the register windowing technique, so that the software can only access 16 of the 64 AR registers at any given time. The programming performance can be effectively improved by rotating windows, replacing function calls, and saving registers when exceptions are triggered.
\end{TableNotes}
\begin{longtable}[c]{|l|c|c|c|l|}
    \caption{Register List of \chipname{} Extended Instruction Set}
    \label{tab:pie-regs-list}\\\hline
    \rowcolor{lightgray}
    \textbf{Register Mnemonics} & \textbf{Quantity} & \textbf{Bit Width} & \textbf{Access} & \textbf{Type} \\\hline
    \endfirsthead
    \multicolumn{4}{c}{\bfseries \tablename\ \thetable{} -- cont'd from previous page} \\\hline
    \rowcolor{lightgray}
    \textbf{Register Mnemonics} & \textbf{Quantity} & \textbf{Bit Width} & \textbf{Access} & \textbf{Type} \\\hline
    \endhead
    
    \multicolumn{4}{r}{\textbf{Cont'd on next page}} \\
    \endfoot
    \insertTableNotes
    \endlastfoot
    AR & 16\tnote{1} & 32 & R/W & General-purpose registers \\\hline
    FR & 16 & 32 & R/W & General-purpose registers to FPU \\\hline
    QR & 8 & 128 & R/W & Customized general-purpose registers \\\hline
    SAR & 1 & 6 & R/W & Special register \\\hline
    SAR\_BYTE & 1 & 4 & R/W & Customized special register \\\hline
    ACCX & 1 & 40 & R/W & Customized special register \\\hline
    QACC\_H & 1 & 160 & R/W & Customized special register \\\hline
    QACC\_L & 1 & 160 & R/W & Customized special register \\\hline
    FFT\_BIT\_WIDTH & 1 & 4 & R/W & Customized special register \\\hline
    UA\_STATE & 1 & 128 & R/W & Customized special register \\\hline
\end{longtable}


\end{ThreePartTable}